\section{Scope, logical foundation, and the meaning of a theorem}
\label{sec:scope}

The aim is a useful foundation for continuous probability inside
intuitionistic primitive recursive arithmetic, abbreviated iPRA. The main
outcome is an explicit arithmetic interpretation of a substantial represented
fragment: continuous and singular continuous distributions, integration,
products, conditioning with separation data, quantitative convergence, and
laws of large numbers and a central limit theorem. It is not an axiom asserting
that the full classical universe of measurable sets exists inside a
first-order number theory.

The subject outline in the requested probability-theory article motivates
the choice of topics \cite{wikiProbability}. Constructive probability has a
large existing literature; Chan's systematic treatment, for example, develops
probability and stochastic processes in Bishop-style constructive analysis
\cite{chan}. Our additional restriction is proof-theoretic: each computation
must have a primitive recursive description, and the induction used in the
object theory must be quantifier-free. A constructive argument in a stronger
analysis is not automatically an iPRA argument. The following presentation
makes the finite calculations and the quantitative hypotheses visible.

\subsection{The exact object theory}

Use first-order intuitionistic natural deduction with equality, the
connectives $\bot,\wedge,\vee,\to$, and the quantifiers $\forall,\exists$.
Negation abbreviates implication to $\bot$. The nonlogical language has
$0,S$ and a function symbol for each \emph{fixed primitive recursive
definition}. Its axioms are the defining equations, $Sx\ne0$, injectivity
of $S$, and quantifier-free induction:
\[
 A(0,\mathbf u)\wedge
 \forall n\bigl(A(n,\mathbf u)\to A(Sn,\mathbf u)\bigr)
 \ \longrightarrow\ \forall n\,A(n,\mathbf u).
\]
Here $A$ is quantifier-free. A recursion symbol has equations
\[
 f(0,\mathbf u)=g(\mathbf u),\qquad
 f(Sn,\mathbf u)=h(n,f(n,\mathbf u),\mathbf u).
\]
The logical inference rules are the usual intuitionistic introduction and
elimination rules, equality substitution, and weakening/substitution of
derivations. No excluded-middle axiom, Markov principle, countable choice,
or comprehension principle is added. Decidable finite rational comparisons
are justified arithmetic case splits.

This is the theory implemented by ProveIt's
\texttt{ConciseIPRA/PR.lean} and its natural-deduction library
\cite{repoCore}. A typed Lean proof of a theorem about the standard natural
numbers is a different object from a term of that embedded derivation type.
The report keeps these two verification levels separate.

\subsection{Numerical coding and total finite operations}

A signed integer is represented by a pair $(u,v)$ with value $u-v$;
cross multiplication gives equality and order without a quotient operation.
A rational code is $(u,v,d)$ with denominator $d+1$, denoting
$(u-v)/(d+1)$. Redundant codes are harmless. Addition, multiplication,
absolute value, comparison, min/max, and division by a denominator with a
supplied positive lower bound are fixed PR operations on these codes.
For malformed optional data use a declared default value; the validity
predicate guards each theorem. This keeps the underlying operations total.

Pairs and finite lists have standard bounded-search codes. Concretely a
finite list of natural entries can be stored as $(n,B,c)$ with $B\ge2$,
$c=\sum_{i<n}a_iB^i$, and $a_i<B$. Entry extraction is
$\lfloor c/B^i\rfloor\bmod B$. For a freshly constructed list take $B$
larger than its finite maximum. Length, concatenation, sorting by bounded
insertion, prefix sums, products, and table refinement are PR. Exponentially
large tables remain primitive recursive; no claim of computational
efficiency follows just from this classification.

Every bounded Boolean condition has an explicit characteristic function.
For example, a universal row condition becomes a fold of conjunctions over
the coded row count. Its connection with the displayed bounded formula is
proved by quantifier-free induction on that fold. A proof must establish this
reflection statement; running a checker on examples is insufficient.

\subsection{How continuous notation is translated}

The phrase ``let $X$ be a represented random variable'' fixes PR programs for
its approximants, their finite levels, and the required error schedules,
together with proofs of their stated contracts. Natural parameters can remain
free. The statement is a schema over these fixed definitions, or a theorem
over a separately specified bounded program interpreter with a proved clock.
It is not quantification in iPRA over arbitrary functions $\mathbb N\to\mathbb N$.

For example, the regularity of an event name $A_k$ with levels $\ell(k)$ is
the first-order formula
\[
 \forall k,l\quad
 2^{k+l}\operatorname{DiffCount}(A_k,A_l)
 \le 2^{\max(\ell(k),\ell(l))}(2^k+2^l).
\]
The count compares the two tables on their common refinement. Its matrix
is an integer inequality. A finite proof of an output estimate asks for
finitely many instances of the input contracts, uses finite arithmetic, and
then universally generalizes the free precision. One must not perform
induction on the all-precision assertion itself.

Here is an explicit denominator-cleared probability theorem. Let $w_i,a_i$
be natural weights and nonnegative numerators, $W=\sum_{i<n}w_i>0$, and
$D,b>0$. Markov's inequality for values $a_i/D$ and threshold $c/b>0$ is
the single finite inequality
\begin{equation}\label{eq:markov-object}
 cD\sum_{i<n}w_i[ba_i\ge cD]
 \le b\sum_{i<n}w_i a_i.
\end{equation}
The partial-sum version of this same inequality is a quantifier-free
induction invariant. The displayed indicator is a PR Boolean comparison.
After proof, its usual probability interpretation divides by $bDW$.

Similarly, for signed $a_i$, a center $m/(WD)$, and threshold $c/b$,
Chebyshev's power-Markov kernel is
\[
 (cWD)^2\sum_{i<n}w_i[b|Wa_i-m|\ge cWD]
 \le b^2\sum_{i<n}w_i|Wa_i-m|^2.
\]
Only when $m=\sum_iw_ia_i$ has been established is the center the mean.
The continuous results preserve this discipline: all analytic error terms
are replaced by rational allowances before the final arithmetic step.

\subsection{The relation to existing ProveIt work}

Two source snapshots were inspected. The local arithmetic and probability
providers are pinned to
\texttt{b2341af5370cd8a43f72ebb2341c0d7d7395b634}. The newer measure and
observation reports were fetched at the then-current main commit
\texttt{83421593771f877bc1ae98f4229eb35507bff05b}.
The repository's directory README distinguishes the checked
\texttt{ConciseIPRA}, \texttt{Derivations}, and \texttt{Bridge} libraries
from research packets retained as delivered. Several newer packets report
their own checked finite Lean kernels, which still do not amount to full
embedded analytic derivations. No repository-wide rebuild or audit is claimed.

This work reuses the coherent-partition and completion approach, excess-tail
clipping, and finite observation discipline from those newer reports
\cite{repoObservable,repoMeasureTransport}. The existing expectation,
normal-distribution, and CLT packets provide useful thematic context
\cite{repoProbability}. Their finite algebra is partly implemented by actual
\texttt{Prf PRAAx} providers. The new package connects these ingredients into
a continuous probability foundation, gives an explicit coherent density
realization, and supplies a self-contained conservative CLT proof. It makes
no claim to have invented the underlying classical or constructive theorems.

\section{The finite analysis needed for probability}
\label{sec:finite-analysis}

The analytic input is small and quantitative. This section specifies the
finite content of integration and calculus used later, so that a reference
to an integral does not silently assume a classical measure construction.

\subsection{Compact quadrature as a Cauchy construction}

A continuous-function presentation on a rational interval $[a,b]$ supplies
a rational evaluator $\widehat f(x,p)$ at rational $x$, with error at most
$2^{-p}$, a rational bound $M$, and a PR modulus $\omega$ satisfying
\[
 |x-y|\le2^{-\omega(k)}\ \Longrightarrow\ |f(x)-f(y)|\le2^{-k}.
\]
These are contracts on regular real names, with rational guards in the
finite implementation. Piecewise presentations carry a finite list of
breakpoints and one-sided values. The integrals of the functions considered
here are constructed only from such data or from explicitly stronger norm
approximations.

For a partition of mesh at most $2^{-\omega(k)}$, tagged rational quadrature
with evaluation error $2^{-p}$ has uncertainty at most
\[
 (b-a)(2^{-k}+2^{-p}).
\]
To prove this, refine two proposed partitions to their finite common list of
endpoints. Each tag differs from any point in its cell by the mesh bound, so
the finite upper and lower rectangle bounds enclose its value with the
displayed error radius. Their symmetric enclosure has twice that radius;
comparing two quadratures adds their error radii. Multiplying by cell lengths
and summing bounds the difference between sufficiently fine quadratures. It produces a
regular real name for the integral. No supremum or infimum over arbitrary
functions is used. A Lipschitz constant $L$ gives the simpler mesh allowance
$L(b-a)h$.

All continuous operations used below come with local computable bounds.
The integral is positive, additive over adjacent intervals, and linear,
because these statements hold for the common finite sums and their errors
can be made arbitrarily small. An improper integral is admitted only when
its finite annular tails have a PR Cauchy bound. In particular the hypothesis
(K1) below uses finite intervals and is not a circular assumption about an
already-existing infinite integral.

\subsection{Finite calculus identities}

The form of the fundamental theorem used here is the following. Suppose
$g,g'$ have the above presentations and a proved local increment estimate
\[
 |g(x+h)-g(x)-hg'(x)|\le Ah^2/2
\]
on a compact interval. Sum it over a partition of mesh $h_0$. The left
increments telescope, so
\[
 \left|g(b)-g(a)-\sum_i g'(x_i)(x_{i+1}-x_i)\right|
 \le A(b-a)h_0/2.
\]
The explicit vanishing error proves $\int_a^b g'=g(b)-g(a)$.
More generally one may supply a proved increment contract
$|g(x+h)-g(x)-hg'(x)|\le |h|\rho(|h|)$ with a PR decay modulus for $\rho$;
the error is then $(b-a)\rho(h_0)$. A bare continuity assertion about a
derivative is not substituted for this quantitative contract. Integration by parts follows by applying
the same finite increment calculation to the product $fg$ and using its
derivative identity. The exponential increment and derivative estimates
follow from its finite power series and geometric remainders; rational and
polynomial functions use exact divided differences and denominator margins.
Thus the applications to normal and exponential kernels require no general
differentiation oracle.

The one-dimensional substitutions used later also have a finite proof.
For a monotone differentiable map $\psi$ on $[a,b]$, carry a bound for
$|\psi'|$, a derivative modulus, and the local modulus of $f$ on its image.
Subdivide the domain. Replace each image-cell length
$\psi(x_{i+1})-\psi(x_i)$ by
$\psi'(x_i)(x_{i+1}-x_i)$ using the increment bound. Oscillation of $f$
on each image cell controls its tagged-sum error. The resulting explicit
error tends to zero and gives
$\int_{\psi(a)}^{\psi(b)}f=\int_a^b(f\circ\psi)\psi'$.
Approximating represented endpoints contributes at most the endpoint error
times a local integrand bound. In this report the maps are affine maps or
reciprocal maps on intervals separated from zero, where every such bound
is a displayed rational expression. For decreasing maps reverse the
endpoints first, or retain oriented interval lengths throughout the finite
sum identity. Improper substitutions separately carry
their spatial-tail estimates.

On compact rectangles, Fubini is equality of finite double sums followed by
the same oscillation bounds. For an unbounded product integral we first
truncate so the \emph{absolute} omitted integral has a supplied bound. This
is the finite Fubini fragment used for products and Gaussian normalization.
It does not make the value of an arbitrary $L^1$ equivalence class on one
specified measure-zero section meaningful.

For the numerical normal example in the archive, the composite midpoint
error is particularly convenient. A twice differentiable function with
$|f''|\le A$ satisfies
\[
 \left|\int_a^b f-h\sum_i f(x_i+h/2)\right|
 \le A(b-a)h^2/24.
\]
Integrate its second-order pointwise remainder symmetrically on each cell;
the linear term cancels, and $\int_{-h/2}^{h/2}t^2/2\,dt=h^3/24$.
For $w(x)=e^{-x^2/2}$,
$|w''(x)|\le (1+x^2)e^{-x^2/2}\le3$, using $te^{-t}\le1$.
This explains the exact enclosure calculation rather than treating numerical
agreement with a floating-point library as evidence of proof.
